\documentclass[12pt,a4paper]{article}

\usepackage[utf8]{inputenc}
\usepackage[T2A]{fontenc}
\usepackage[russian]{babel}
\usepackage[margin=2cm]{geometry}
\usepackage{listings}
\usepackage{xcolor}
\usepackage{booktabs}
\usepackage[hidelinks]{hyperref}
\usepackage{graphicx}

\lstset{
    basicstyle=\ttfamily\small,
    breaklines=true,
    frame=single,
    columns=fullflexible,
    backgroundcolor=\color{gray!5}
}

\begin{document}

\begin{titlepage}
\begin{center}
    \includegraphics[width=\textwidth]{mai_header.jpg}

    \vspace{1cm}
    Министерство науки и высшего образования Российской Федерации\\
    Федеральное государственное бюджетное образовательное учреждение высшего образования\\
    \textbf{Московский авиационный институт}\\
    (Национальный исследовательский университет)

    \vspace{1.5cm}
    Институт № 8 «Компьютерные науки и прикладная математика»\\
    Кафедра 806

    \vspace{2cm}
    {\large\textbf{Отчёт по лабораторной работе №1}}\\[0.3cm]
    {\textbf{Дисциплина}}: «Операционные системы и архитектура компьютера»\\[0.3cm]
    {\textbf{Тема}}: «Анализ профиля системных вызовов программы \texttt{ffmpeg}»\\[0.3cm]
\end{center}

\vspace{2cm}
\begin{flushright}
    {\textbf{Выполнил:}}\\ 
    студент группы М8О-206БВ-25\\
    Ширанков М.С.\\

    {\textbf{Проверил:}}\\
    Ассистент Центра ТОП-ИТ\\
    Ефимов Д.А.
\end{flushright}

\vfill
\begin{center}
    Москва 2026
\end{center}
\end{titlepage}

\tableofcontents
\newpage

\section{Цель работы}
Построить профиль системных вызовов программы (\texttt{ffmpeg})
с помощью утилиты \texttt{strace}, получить статистику вызовов и
разбить её по логическим подсистемам ядра (память, процессы/потоки,
файловый ввод-вывод, сигналы, сеть), не смешивая их между собой.

\section{Ход работы}
В качестве исследуемой команды используется перекодирование видео
\texttt{zazu\_input.mp4} с явным указанием видеокодека
\texttt{libx264}:

\begin{lstlisting}
ffmpeg -i zazu_input.mp4 -c:v libx264 zazu_output.mp4
\end{lstlisting}

Команда задействует сразу несколько подсистем ОС:
\begin{itemize}
    \item чтение исходного видео и запись результирующего файла (файловый I/O);
    \item динамическую загрузку разделяемых библиотек (\texttt{libavcodec},
          \texttt{libx264} и др.) при старте процесса;
    \item многопоточное кодирование кадров библиотекой \texttt{libx264}
          (создание потоков и их синхронизацию);
    \item управление памятью под буферы кадров и рабочие структуры кодека.
\end{itemize}

Ниже разобран построчный (не агрегированный) лог одного запуска
\texttt{strace -f -s 200 -tt ffmpeg -i zazu\_input.mp4 -c:v libx264 out.mp4}
по пяти хронологическим фазам — от старта процесса до завершения.

\subsection{Запуск процесса}
Первая группа вызовов — старт процесса и работа динамического загрузчика:

\begin{lstlisting}
execve("/usr/bin/ffmpeg", ["ffmpeg", "-i", "zazu_input.mp4", ...]) = 0
access("/etc/ld.so.preload", R_OK) = -1 ENOENT
openat(AT_FDCWD, "/etc/ld.so.cache", O_RDONLY|O_CLOEXEC) = 3
fstat(3, ...) = 0
mmap(NULL, 70655, PROT_READ, MAP_PRIVATE, 3, 0) = 0x...
close(3) = 0
openat(AT_FDCWD, "/lib/x86_64-linux-gnu/libavdevice.so.60", ...) = 3
read(3, "\177ELF...", 832) = 832
fstat(3, ...) = 0
mmap(...)  mmap(...)  mmap(...)
close(3) = 0
\end{lstlisting}

Эта пятёрка \texttt{openat → read → fstat → mmap(×2-4) → close} повторяется
для каждой из более чем 100 загружаемых разделяемых библиотек: сначала
основные компоненты \texttt{ffmpeg} (\texttt{libavdevice}, \texttt{libavfilter},
\texttt{libavformat}, \texttt{libavcodec}, \texttt{libavutil}, \texttt{libx264.so.164}),
затем их транзитивные зависимости. После загрузки последней библиотеки
главный поток инициализируется:

\begin{lstlisting}
arch_prctl(ARCH_SET_FS, 0x...) = 0
set_tid_address(0x...) = 29257
set_robust_list(0x..., 24) = 0
rseq(0x..., 0x20, 0, 0x53053053) = 0
mprotect(0x..., 16384, PROT_READ) = 0
mprotect(0x..., 4096, PROT_READ) = 0
...
\end{lstlisting}
Далее следует ещё по одному \texttt{mprotect} на каждый сегмент каждой
загруженной библиотеки.

Серия \texttt{mprotect(..., PROT\_READ)} переводит только что отображённые
сегменты библиотек из \texttt{writable} (нужного на время линковки) в
\texttt{read-only} — запись в них больше не требуется.

\subsection{Настройка окружения}
\begin{lstlisting}
sched_getaffinity(0, 512, [0 1 2 3]) = 16
openat(AT_FDCWD, "zazu_input.mp4", O_RDONLY) = 3
write(2, "Input #0, mov,mp4,m4a,3gp,3g2,mj2, from 'zazu_input.mp4':\n", 58) = 58
\end{lstlisting}
\texttt{sched\_getaffinity} сообщает ffmpeg доступные ядра CPU (здесь — 4,
\texttt{[0 1 2 3]}); это число определяет размер будущего пула потоков
кодирования. Затем открывается и пробуется (probing) входной контейнер —
результат распознавания формата (\texttt{mov,mp4,...}) ffmpeg печатает в stderr.

\subsection{Построение конвейера}
У \texttt{ffmpeg} кодек \texttt{libx264} слинкован в тот же процесс, поэтому
"конвейер" обработки — это не ОС-пайп, а граф \textbf{потоков} с общей памятью
(\texttt{clone3} с флагом \texttt{CLONE\_VM}), синхронизируемых через
\texttt{futex} вместо \texttt{read}/\texttt{write} по дескриптору:

\begin{lstlisting}
openat(AT_FDCWD, "out.mp4", O_WRONLY|O_CREAT|O_TRUNC, 0666) = 4
clone3({flags=CLONE_VM|CLONE_FS|CLONE_FILES|CLONE_SIGHAND|CLONE_THREAD|...,
        stack=0x..., stack_size=0x7ff400, ...}) = 29304
prctl(PR_SET_NAME, "av:h264:df0") = 0
...
prctl(PR_SET_NAME, "av:h264:df4") = 0
prctl(PR_SET_NAME, "dec0:0:h264") = 0
prctl(PR_SET_NAME, "dmx0:mov,mp4,m4a") = 0
...
prctl(PR_SET_NAME, "mux0:mp4") = 0
\end{lstlisting}

Из 19 вызовов \texttt{clone3} восемь потоков получают явное имя через
\texttt{prctl(PR\_SET\_NAME, ...)}, которое прямо видно в трассировке и
раскрывает внутреннюю схему конвейера: \texttt{dmx0} — поток демультиплексора
(разбирает контейнер на потоки), \texttt{dec0:0:h264} — поток декодера
видеодорожки, \texttt{av:h264:df0}…\texttt{df4} — пять потоков покадрового
декодирования H.264 (frame-threading), \texttt{mux0} — поток мультиплексора,
формирующий выходной MP4. Оставшиеся \textasciitilde11 потоков из 19 —
безымянный пул worker-потоков кодера \texttt{libx264} (slice-based threading);
каждый получает собственный стек \textasciitilde8 МБ через
\texttt{mmap(PROT\_NONE)}+\texttt{mprotect}.

\subsection{Кодирование и запись данных}
Основная масса вызовов в этой фазе — \texttt{futex}: декодер передаёт готовые
кадры энкодеру, энкодер раздаёт срезы worker-потокам, каждый из них
сигнализирует о готовности через \texttt{FUTEX\_WAKE}/\texttt{FUTEX\_WAIT\_BITSET} —
именно эта активность и даёт те самые 98\% времени в системных вызовах
из раздела~\ref{sec:categories}.

В выходной файл (дескриптор 4) производится всего \textbf{4 вызова}
\texttt{write}:
\begin{lstlisting}
write(4, "\0\0\0 ftypisom...free...mdat", 48) = 48
write(4, "...x264 - core 164 ... options: cabac=1 ref=3...", 195063) = 195063
write(4, "\0\2\371\377", 4) = 4
write(4, "\0\0\6\203moov\0\0\0lmvhd...", 1667) = 1667
\end{lstlisting}
Первый вызов — заголовок контейнера (\texttt{ftyp}/\texttt{free}/заглушка
\texttt{mdat}); второй — одним вызовом весь закодированный битовый поток H.264
(195\,063 байта, с различимым SEI-префиксом x264 с версией кодека и
параметрами кодирования); третий — служебные байты; четвёртый —
атом \texttt{moov} (индекс фильма), который MP4-мультиплексор по умолчанию
дописывает в конец файла только после того, как все кадры уже записаны.

\subsection{Завершение процесса}
Когда кодирование окончено, главный поток завершает worker-потоки один за
другим — это ядро реализации \texttt{pthread\_join}, построенное на \texttt{futex}:
\begin{lstlisting}
futex(0x..., FUTEX_WAKE_PRIVATE, 1) = 1
rt_sigprocmask(SIG_BLOCK, ~[RT_1], NULL, 8) = 0
madvise(0x..., 8368128, MADV_DONTNEED) = 0
exit(0) = ?
+++ exited with 0 +++
futex(0x..., FUTEX_WAIT_BITSET|FUTEX_CLOCK_REALTIME, ..., NULL, ...) = 0
munmap(0x..., 8392704) = 0
\end{lstlisting}
Главный поток будит очередной рабочий поток (\texttt{FUTEX\_WAKE\_PRIVATE});
тот блокирует сигналы, освобождает свой стек (\texttt{madvise(MADV\_DONTNEED)})
и вызывает \texttt{exit(0)} — завершается только сам поток, а не весь процесс
(в отличие от \texttt{exit\_group}). Главный поток просыпается из
\texttt{FUTEX\_WAIT\_BITSET} (ожидание на \texttt{child\_tid}, выставленного
флагом \texttt{CLONE\_CHILD\_CLEARTID} при создании потока) и освобождает его
стек через \texttt{munmap}. Это повторяется по разу на каждый из 19 потоков.
После этого главный поток закрывает выходной файл
(\texttt{close(4)}) и завершает процесс целиком: \texttt{exit\_group(0)}.

\section{Получение полного профиля вызовов}

Флаг \texttt{-c} переводит \texttt{strace} в режим сбора статистики: вместо построчной
трассировки каждого вызова выводится статистическая таблица (число вызовов,
суммарное и среднее время, количество вызовов, завершившихся ошибкой). Флаг \texttt{-f}
обязателен для \texttt{ffmpeg}, так как кодек \texttt{libx264} создаёт дочерние потоки
через \texttt{clone3()} - без \texttt{-f} \texttt{strace} видит только главный поток и
теряет всю статистику по кодированию.

\begin{lstlisting}
strace -f -c ffmpeg -i zazu_input.mp4 -c:v libx264 zazu_output.mp4
\end{lstlisting}

Результат выполнения (перекодирование тестового видео \texttt{zazu\_input.mp4} на
4-ядерной машине):

\begin{lstlisting}
% time     seconds  usecs/call     calls    errors syscall
------ ----------- ----------- --------- --------- ------------------
 98.19    7.442150        1842      4039       330 futex
  0.63    0.047605          44      1063           mmap
  0.29    0.021859          51       421           mprotect
  0.27    0.020194         196       103           munmap
  0.12    0.009123          37       242           read
  0.12    0.008736          38       228         5 openat
  0.09    0.006535          28       229           fstat
  0.08    0.005932          26       223           close
  0.06    0.004730          32       147           write
  0.04    0.003178         167        19           clone3
  0.03    0.002501          32        77           rt_sigprocmask
  0.03    0.002185          66        33           brk
  0.01    0.000900          45        20           rseq
  0.01    0.000861          45        19           madvise
  0.01    0.000823          41        20           set_robust_list
  0.01    0.000426          30        14         4 prctl
  0.00    0.000195          27         7           lseek
  0.00    0.000163          27         6           pselect6
  0.00    0.000155          31         5           sched_getaffinity
  0.00    0.000129          64         2         2 ioctl
  0.00    0.000116          19         6           rt_sigaction
  0.00    0.000116          38         3         3 access
  ...
------ ----------- ----------- --------- --------- ------------------
100.00    7.579191        1090      6948       348 total
\end{lstlisting}

\subsection{Значение столбцов}
\begin{itemize}
    \item \textbf{\% time} - доля суммарного времени выполнения системных вызовов,
          пришедшаяся на данный вызов (не доля времени всей программы - I/O-ожидание
          и вычисления в пространстве пользователя сюда не входят);
    \item \textbf{seconds} - суммарное время, потраченное ядром на обработку всех
          вызовов этого типа;
    \item \textbf{usecs/call} - среднее время одного вызова в микросекундах;
    \item \textbf{calls} - сколько раз вызов был выполнен (суммарно по всем потокам,
          так как использован \texttt{-f});
    \item \textbf{errors} - сколько из этих вызовов завершились с кодом ошибки
          (\texttt{-1} и установленным \texttt{errno});
    \item \textbf{syscall} - имя системного вызова.
\end{itemize}

\subsection{Анализ результата}
Основная долю времени, проведённого в системных вызовах, приходится на \texttt{futex}.
\texttt{Libx264} распараллеливает кодирование кадров/срезов (slice-based
threading) между рабочими потоками по числу ядер CPU (здесь - 4), а сами потоки
создаются вызовом \texttt{clone3}. Синхронизация между ними (очередь заданий, условные переменные, мьютексы
в реализации NPTL/glibc) реализована поверх futex, поэтому именно он доминирует
и по числу вызовов, и по суммарному времени. Остальные вызовы (\texttt{mmap},
\texttt{mprotect}, \texttt{munmap}, \texttt{brk}, \texttt{madvise}) - это управление
памятью под буферы кадров и стеки потоков, а \texttt{read/openat/fstat/close/write} -
чтение исходного файла, загрузка динамических библиотек кодеков при старте и запись
результата.


\section{Разделение профиля по логическим категориям}
\label{sec:categories}

Готовые классы \texttt{strace} (\texttt{\%memory}, \texttt{\%process}, \texttt{\%signal},
\texttt{\%network}, \texttt{\%desc}, \texttt{\%file}, \texttt{\%ipc}, \texttt{\%clock})
позволяют собрать статистику только по вызовам одной подсистемы, не смешивая их с
остальными. Для каждой категории трассировка запускалась отдельным прогоном, чтобы категории физически не
пересекались в выводе:

\begin{lstlisting}
# memory: mmap/munmap/mprotect/brk/madvise
strace -f -c -e trace=%memory  -o mem_profile.log \
    ffmpeg -i zazu_input.mp4 -c:v libx264 zazu_output.mp4

# processes/threads: clone3/execve/wait4/exit_group
strace -f -c -e trace=%process -o process_profile.log \
    ffmpeg -i zazu_input.mp4 -c:v libx264 zazu_output.mp4

# signals: rt_sigaction/rt_sigprocmask/tgkill
strace -f -c -e trace=%signal  -o signal_profile.log \
    ffmpeg -i zazu_input.mp4 -c:v libx264 zazu_output.mp4

# file descriptors / I/O: read/write/open/close/fstat...
strace -f -c -e trace=%desc    -o io_profile.log \
    ffmpeg -i zazu_input.mp4 -c:v libx264 zazu_output.mp4

# network: socket/connect/send*/recv*
strace -f -c -e trace=%network -o net_profile.log \
    ffmpeg -i zazu_input.mp4 -c:v libx264 zazu_output.mp4
\end{lstlisting}

\subsection{Категория: память (\texttt{\%memory})}
\begin{lstlisting}
% time     seconds  usecs/call     calls    errors syscall
------ ----------- ----------- --------- --------- ----------------
 51.55    0.046484          43      1063           mmap
 23.83    0.021487          51       421           mprotect
 20.51    0.018489         179       103           munmap
  3.16    0.002851         150        19           madvise
  0.86    0.000774          23        33           brk
  0.09    0.000079          79         1           mlock
------ ----------- ----------- --------- --------- ----------------
100.00    0.090164          54      1640           total
\end{lstlisting}
Преобладают \texttt{mmap}/\texttt{mprotect}/\texttt{munmap} - это загрузка
динамических библиотек кодеков (\texttt{ld.so} отображает каждую \texttt{.so} в
память и выставляет права страниц) и выделение буферов под кадры или потоковые стеки.

\subsection{Категория: процессы и потоки (\texttt{\%process})}
\begin{lstlisting}
% time     seconds  usecs/call     calls    errors syscall
------ ----------- ----------- --------- --------- ----------------
 82.63    0.002560         134        19           clone3
 17.37    0.000538         538         1           execve
------ ----------- ----------- --------- --------- ----------------
100.00    0.003098         154        20           total
\end{lstlisting}
Один \texttt{execve} - запуск самого процесса \texttt{ffmpeg}; 19 вызовов
\texttt{clone3} - создание потоков пула кодирования \texttt{libx264}.
При просмотре построчного лога с \texttt{-f} по такому
количеству потоков часто встречаются строки вида \texttt{<unfinished ...>} и
\texttt{<... resumed>}. \texttt{Strace} показывает чередование: если, пока один поток ожидает завершения системного вызова
(например, блокирующего \texttt{futex} или \texttt{read}), планировщик
переключается на другой поток и тот тоже входит в syscall, \texttt{strace}
не может напечатать оба вызова как непрерывные строки - начатый, но ещё не
завершившийся вызов помечается \texttt{<unfinished ...>}, а когда управление
возвращается к его поток и вызов действительно завершается, соответствующая
строка получает пометку \texttt{<... resumed>} с итоговым результатом.
 
\subsection{Категория: сигналы (\texttt{\%signal})}
\begin{lstlisting}
% time     seconds  usecs/call     calls    errors syscall
------ ----------- ----------- --------- --------- ----------------
 93.84    0.002454          31        77           rt_sigprocmask
  6.16    0.000161          26         6           rt_sigaction
------ ----------- ----------- --------- --------- ----------------
100.00    0.002615          31        83           total
\end{lstlisting}
\texttt{rt\_sigprocmask} вызывается при каждом создании потока (glibc маскирует
сигналы на время инициализации нового потока), \texttt{rt\_sigaction} - установка
обработчиков (например, \texttt{SIGINT} для корректного завершения по Ctrl+C).

\subsection{Категория: файловый ввод-вывод (\texttt{\%desc})}
\begin{lstlisting}
% time     seconds  usecs/call     calls    errors syscall
------ ----------- ----------- --------- --------- ----------------
 57.66    0.046460          43      1063           mmap
 10.93    0.008807          38       228         5 openat
  9.92    0.007994          32       243           read
  8.11    0.006533          28       229           fstat
  7.93    0.006390          28       223           close
  4.39    0.003540          24       147           write
  0.39    0.000312          44         7           pselect6
  0.26    0.000209          29         7           lseek
------ ----------- ----------- --------- --------- ----------------
100.00    0.080570          37      2158         9 total
\end{lstlisting}
Здесь видно чтение исходного \texttt{zazu\_input.mp4}, запись результата и 
множество пар \texttt{openat}/\\\texttt{fstat}/\texttt{close} - это динамический
загрузчик ищет и открывает разделяемые библиотеки кодеков при старте процесса
(5 из 228 \texttt{openat} завершились ошибкой \texttt{ENOENT} - обычный перебор
путей поиска \texttt{.so}-файлов).

\subsection{Категория: сеть (\texttt{\%network})}
Вывод пуст. Поскольку и вход, и выход - локальные файлы, \texttt{ffmpeg} в этом
сценарии ни разу не обращается к сетевому стеку (ни \texttt{socket}, ни
\texttt{connect}, ни \texttt{sendto/recvfrom}).

\section{Итог работы}
В результате выполнения лабораторной работы были получены практические
навыки профилирования реальной многопоточной программы: сбор статистики (\texttt{-c}),
разделение по классам системных вызовов (\texttt{-e trace=\%class}),
а также разбор сырых трасс для глубокого изучения
понимания работы программы.

\end{document}
